\section{তথ্য ও যোগাযোগ প্রযুক্তি এবং অর্থনৈতিক উন্নয়ন}
\label{sec:ch1-economy}
\subsection{দক্ষতা, বাজার ও সেবা}
ICT উৎপাদন, সেবা, বিপণন, অর্থ লেনদেন, প্রশাসন ও কাজের তথ্য দ্রুত আদান-প্রদান করে। তথ্যের ঘাটতি, যোগাযোগের ব্যয় ও বাজারের দূরত্ব কমলে পরিকল্পনা ও দক্ষতা উন্নত হতে পারে। শিক্ষা দক্ষতা বাড়ায়; স্বাস্থ্যসেবা কাজের সক্ষমতায় সহায়তা করে; কৃষির আবহাওয়া, বাজারদর ও ফসলের তথ্য সিদ্ধান্তে সাহায্য করে।
সরকারি আবেদন, কর ও সনদ অনলাইনে দেওয়া প্রশাসনিক ব্যয় ও সময় কমাতে পারে। মোবাইল আর্থিক সেবায় দ্রুত অর্থ পাঠানো এবং স্থানীয় ব্যবসায় লেনদেন সহজ হয়। অনলাইন দোকান দূরের ক্রেতায় পৌঁছায়। এসব সেবার সংজ্ঞা ও কাজ \ref{sec:ch1-components} অংশে একত্রে আছে।
\chOnePair{\chOneNative{Economy}}{দক্ষতা, নিরাপত্তা, অবকাঠামো ও ন্যায্য প্রবেশাধিকারের সঙ্গে প্রযুক্তি ব্যবহার করলে বাজার ও সেবার সুযোগ বাড়ে।}
\subsection{কর্মসংস্থান ও উদ্যোগ}
দক্ষ ICT জনশক্তি সফটওয়্যার ও সেবা উৎপাদন, বৈধ রপ্তানি আয় এবং উদ্যোগে ভূমিকা রাখে। গ্রামের তরুণ অনলাইন বাজারে গ্রাফিক্স বা ওয়েব উন্নয়নের কাজ করে দক্ষতা, স্থানীয় কর্মসংস্থান ও বৈদেশিক আয়ের সুযোগ তৈরি করতে পারেন। এই আয় কাজ পাওয়া ও যথাযথভাবে শেষ করার ওপর নির্ভর করে।
Freelancing স্বাধীন কর্মসম্পর্ক; Remote Job কাজের স্থান বোঝায়, যেখানে নিয়মিত চাকরিও হতে পারে। Outsourcing-এর সংজ্ঞা \ref{sec:ch1-outsourcing} অংশে আছে। বিদেশে কাজ দেওয়া হলে Offshoring-এর সম্পর্ক থাকতে পারে; সব Outsourcing বিদেশে হয় না।
\subsection{আয়, লাভ ও উপযোগিতা}
স্থানীয় কারিগরের পণ্যের বাজার বাড়াতে অনলাইন দোকানে মান, স্পষ্ট ছবি, দাম, আদেশ, বৈধ অর্থ পরিশোধ, পৌঁছানো, ফেরতের নিয়ম ও ক্রেতা সহায়তা একসঙ্গে সাজাতে হয়। শুধু ওয়েবসাইট বিক্রি বা লাভ নিশ্চিত করে না।
\term{Revenue} বিক্রির আয়; \term{Profit} নির্ণয়ে পণ্য, সরবরাহ, বিজ্ঞাপন, অর্থ পরিশোধ ও ফেরতের ব্যয় বাদ দিতে হয়। দক্ষতার ঘাটতি, প্রতারণা বা তথ্য হারালে লাভ কমতে পারে। প্রশিক্ষণ, মান যাচাই, নিরাপত্তা ও বিকল্প ব্যবস্থা তাই অর্থনৈতিক উন্নয়নের অংশ।
বাংলাদেশের ডিজিটাল সেবা, শিক্ষা, স্বাস্থ্য, কৃষি, মোবাইল লেনদেন ও উদ্যোক্তা উন্নয়নের সুযোগ কাজে লাগাতে দক্ষতা, যোগাযোগ, বাজার বোঝা ও দায়িত্বশীলতা দরকার। যন্ত্র, সংযোগ ও সেবা ব্যবহারের বৈষম্য কমানোও জরুরি।
\begin{practicebox}
\begin{enumerate}
\item E-commerce কী?
\item Mobile banking অর্থনীতিতে কীভাবে সহায়তা করে?
\item Outsourcing ও freelancing-এর মধ্যে সম্পর্ক ব্যাখ্যা কর।
\item ICT অর্থনৈতিক উন্নয়নে কীভাবে ভূমিকা রাখে—চারটি point লেখ।
\end{enumerate}
\end{practicebox}

